\scanpage{222}
\mspage{77}

à l'application décroissante $\alpha_0:Q\unsure{-|\to}P$, transformant Sup en Inf. \deleted{L'application}

Ici il \deleted{peut} est commode de considérer $\unsure{\sharp}$ et $\flat$ comme des opérations sur $P$ resp $Q$, par

\[
 p^\sharp=\alpha_0\beta_0(p),\qquad q^\flat=\beta_0\alpha_0(q).
 \tag*{(5.19)}
\]

La relation entre ces endomorphismes, et l'opération de $=$ nom sur $\unsure{\mathfrak P|\mathcal P}(\deleted{\unsure{A}}P),\mathfrak P(Q)$, est donnée par

\[
 \left\{\begin{aligned}
 p^\sharp&=\text{\unsure{Plus|plus} grand élément de }\{p\}^\sharp\\
 A^\sharp&=\deleted{\unsure{p}}\mathop{\mathrm{Sup}}_{p\in A}p^\sharp,
 \end{aligned}\right.
 \tag*{(5.20)}
\]

et symétriquement du côté de $Q$.

\deleted{Les applications} On dit que $p$ est \uline{$R$-clos}, ou $\beta_0$-clos, ou $\alpha_0$-clos, si $p=p^\sharp$, et de $=$ $q\in Q$ est dit $R$-\uline{clos} etc si $q=q^\flat$. Les \deleted{appl} applications

\[
 \left\{\begin{aligned}
 p\unsure{\vdash|\mapsto} p^\sharp &: P\longrightarrow P\\
 q\vdash q^\flat &: Q\unsure{-|\to}Q
 \end{aligned}\right.
 \tag*{(5.21)}
\]

sont \uline{croissantes}, \inserted{et \uline{idempotentes},} \deleted{et commutent aux Sup \uline{quelconques}. Ce dernier fait résulte de} \inserted{on peut caractériser} $p^\sharp$ comme le plus petit élément clos de $P$ majorant $p$, et de $=$ pour $q^\flat$. On trouve \unsure{alors}, par (5.19), que $\alpha_0$ et $\beta_0$ induisent des bijections inverses l'une de l'autre,
% archive sheet number: 222
